\documentclass[11pt]{article}
\usepackage[margin=1in]{geometry}
\usepackage{listings}
\usepackage{xcolor}
\usepackage{hyperref}

\lstset{
  basicstyle=\ttfamily\footnotesize,
  breaklines=true,
  frame=single,
  numbers=left,
  numberstyle=\tiny\color{gray},
  keywordstyle=\color{blue},
  stringstyle=\color{red!70!black},
  commentstyle=\color{green!50!black},
  showstringspaces=false,
}

\title{Casino Slots: A Gentle Tour of the Code\\for Web Development Newcomers}
\author{}
\date{}

\begin{document}
\maketitle

\section{Before We Start: What Is This?}

The file \texttt{src/App.jsx} is the heart of a small slot-machine web app.
If you have never built for the web before, here is the mental picture you need:

\begin{itemize}
  \item A \textbf{web page} is just HTML (structure), CSS (style), and JavaScript (behavior) running inside a browser.
  \item \textbf{React} is a JavaScript library that lets you describe what the page should look like as a function of \emph{state} (variables that change over time). When state changes, React re-draws only the parts of the page that need updating.
  \item \textbf{JSX} is the funny-looking syntax where HTML-ish tags (\texttt{<div>}, \texttt{<button>}) appear inside JavaScript. The \texttt{.jsx} file extension signals that the file mixes the two.
\end{itemize}

A React app is made of \textbf{components}: reusable functions that return JSX.
Our file has two components: \texttt{Reel} (one spinning column) and \texttt{App} (the whole machine).

\section{The Top of the File: Imports and Constants}

\begin{lstlisting}[language=JavaScript]
import { useEffect, useMemo, useRef, useState } from 'react';
import logoUrl from '../logo.png';
\end{lstlisting}

\texttt{import} pulls code from elsewhere. The first line grabs four tools from React
called \emph{hooks} (more on these soon). The second line imports an image;
the bundler (Vite, in this project) turns the image into a URL string we can use in an \texttt{<img>} tag.

\begin{lstlisting}[language=JavaScript]
const JACKPOT_SYMBOL = 'LOGO';
const SYMBOLS = [JACKPOT_SYMBOL, 'cherry', 'diamond', ...];
const REEL_COUNT = 3;
const SPIN_DURATION = 2200;
\end{lstlisting}

These are \textbf{constants}: values that never change. Putting ``magic numbers''
like \texttt{2200} (milliseconds) behind named constants makes the code easier to read and tweak.

\section{Pure Helper Functions}

Three plain functions do the math of the game. ``Pure'' means they take
inputs and return outputs without touching anything else.

\begin{description}
  \item[\texttt{buildReels()}] Creates three arrays of 12 symbols each. It uses
  \texttt{Array.from({length: N}, ...)} which is a compact way to say
  ``make an array of N items, computed from the index.''
  \item[\texttt{createVisibleGrid(reelSets, centers)}] Given the current center
  position of each reel, returns the three symbols visible in that reel
  (the one above, the center one, and the one below).
  The \texttt{\%} (modulo) operator wraps the index around so the reel feels infinite.
  \item[\texttt{getLineWin(centerRow)}] Counts how many of each symbol appear on the
  center line. If any symbol shows up at least twice, that is a win.
\end{description}

\section{The \texttt{Reel} Component}

\begin{lstlisting}[language=JavaScript]
function Reel({ symbols, centerIndex, spinning, highlighted }) { ... }
\end{lstlisting}

Components receive data through \textbf{props} (the object inside the braces).
This reel is told which symbols to show, which index is currently centered,
whether it is spinning, and whether it should glow as part of a win.

Inside the function, we build a ``strip'' that repeats the symbols five times
(\texttt{REPEAT\_MULTIPLIER}). We then use a CSS \texttt{transform: translateY(...)}
to slide the strip up so the chosen symbol sits in the window. The spinning
animation itself lives in a CSS file (not shown here) and is switched on by
the class name \texttt{spinning}.

\texttt{useMemo} caches the built strip so React does not recreate it on
every render---a small performance trick.

\section{The \texttt{App} Component: State Hooks}

\begin{lstlisting}[language=JavaScript]
const [credits, setCredits] = useState(STARTING_CREDITS);
\end{lstlisting}

\texttt{useState} is how a component ``remembers'' something between renders.
It returns a pair: the current value, and a function to update it.
Calling \texttt{setCredits(400)} tells React ``the credit count changed, please re-render.''

The app tracks: which symbol each reel is centered on, whether each reel has
settled after a spin, the player's credits, the current bet, a status message,
the last win, and the spin counter.

\section{Refs and Cleanup}

\begin{lstlisting}[language=JavaScript]
const timersRef = useRef([]);
\end{lstlisting}

\texttt{useRef} stores a value that persists across renders but, unlike state,
does \emph{not} trigger a re-render when it changes. We use it to keep
a list of active timers and animation frames so we can cancel them.

\begin{lstlisting}[language=JavaScript]
useEffect(() => {
  return () => { /* cancel timers */ };
}, []);
\end{lstlisting}

\texttt{useEffect} runs code with side effects. The function it \emph{returns}
is a cleanup function. With the empty array \texttt{[]} as the second argument,
the effect runs once when the component mounts, and its cleanup runs once
when the component unmounts. This prevents stray timers from firing after
the app is gone.

\section{Spinning the Reels}

When the user clicks \textbf{Spin}, \texttt{handleSpin()} does the following:

\begin{enumerate}
  \item Refuse to spin if the reels are already moving or credits are too low.
  \item Cancel any lingering timers.
  \item Pick a random final position for each reel.
  \item Deduct the bet from the credits.
  \item Start the animation loop (\texttt{scheduleSpinAnimation}).
  \item Schedule each reel to ``settle'' with a staggered delay (\texttt{STOP\_STAGGER}) so they land one after another.
  \item Schedule \texttt{finishSpin} to run after the last reel lands, which tallies any winnings.
\end{enumerate}

The animation uses \texttt{requestAnimationFrame}, the browser's way of asking
``call me again just before the next repaint.'' Each frame nudges the
visible symbols to new indices, giving the illusion of spinning. When
enough time has passed (\texttt{SPIN\_DURATION}), the reel snaps to its predetermined target.

\section{Rendering: What the User Sees}

The bottom half of \texttt{App} is one big JSX expression.
It describes the layout top-to-bottom:

\begin{itemize}
  \item A \texttt{<main>} wrapper; its class name gains \texttt{win-state} when the player just won, which CSS can use to light things up.
  \item A header with the title and ``Last payout.''
  \item A status deck showing credits, current bet, and spin count.
  \item A grid of three \texttt{<Reel>} components, one per reel.
  \item Controls: bet up/down, max bet, and the big spin button.
\end{itemize}

Notice the patterns:

\begin{itemize}
  \item \texttt{\{credits\}} inside JSX prints the JavaScript value.
  \item \texttt{onClick=\{handleSpin\}} attaches a click handler---no quotes, because it is a function reference, not a string.
  \item \texttt{disabled=\{!canSpin\}} disables the button when spinning is not allowed. \texttt{canSpin} is computed each render from current state.
  \item \texttt{reels.map((reel, i) => <Reel key=\{...\} ... />)} turns an array into a list of components. React requires a unique \texttt{key} on each item so it can tell them apart when the list changes.
\end{itemize}

\section{The One-Way Data Flow}

The whole program follows one rule: \textbf{data flows down, events flow up}.

\begin{itemize}
  \item The \texttt{App} component owns the state.
  \item It passes pieces of that state down to \texttt{Reel} through props.
  \item When the user clicks a button, handlers in \texttt{App} update the state.
  \item React re-renders, the new state flows back down, and the screen updates.
\end{itemize}

Once that loop clicks in your head, most React apps stop looking mysterious.

\section{Where to Poke Around Next}

Good small experiments to build intuition:

\begin{itemize}
  \item Change \texttt{STARTING\_CREDITS} and reload the page.
  \item Add a new emoji to the \texttt{SYMBOLS} array.
  \item Lower \texttt{SPIN\_DURATION} to 800 and feel the difference.
  \item Add a \texttt{<p>} element that shows \texttt{\{spinCount\}} somewhere new.
\end{itemize}

Each change is small, safe, and shows you a different slice of how state,
rendering, and the browser fit together.

\end{document}
